\scanpage{43}
\mspage{130}

On \unsure{voit|note} que cette gerbe $\mathscr G$ est loc. constante
(puisque elle correspond à un lien $f_*(\unsure{G|C})$ loc constant),
et que $\mathscr G'$ est son image inverse. Mais comme
$p$ est une Hot-équivalence, il s'ensuit que
le foncteur
\[
\underline{\Gamma}(Y,\mathscr G)\unsure{\text{---}|\to}\underline{\Gamma}(Y',\mathscr G')
\]
est une équivalence de catégories, donc $\unsure{\xi'|\zeta'}$
provient (à iso. près) d'un unique $\xi$ de $\underline{\Gamma}(Y,\mathscr G)$,
qed.

$3^\circ$) Soit $F$ faisceau abélien loc. constant sur $X$,
il faut voir
\[
H^i(X,F)\xrightarrow{\sim} H^i(X',F')\qquad (F'=q^*F),\tag*{$(*)$}
\]
pour tt $i\in\mathbb N$. Mais \deleted{\unsure{comme}} les deux termes sont
l'aboutissement de suites spectrales de
Leray, de termes initiaux
\[
\begin{aligned}
E_2^{pq}&\simeq H^p(Y,R^q f_*(F))\\
E'{}^{pq}_2&\simeq H^p(Y',R^q f'_*(F')
\end{aligned}
\]
et \unsure{ceci|$(*)$} est associé à un hom de suites
spectrales, qui pour les termes initiaux est
\[
H^p(Y,R^q f_*(F))\unsure{\simeq|\longrightarrow}H^{\unsure{q|p}}(Y',R^q f'_*(F'))\tag*{$(**)$}
\]
associé à
\[
p^*(R^q f_*(F))\text{---}R^q f'_*(F')\tag*{$(***)$}
\]
Mais par hyp. les $R^q f_*(F)$, \inserted{$R^q f'_*(F')$}, sont loc const.
et se calculent fibre par fibre,

% archive sheet number: 43
